% !TeX root = ../../main.tex

PostgreSQL ist eine open-source Datenbank die seit 35 Jahren aktiv entwickelt wird.
Die aktuell veröffentliche Version ist, 18.4.
Version 19 wird gerade entwickelt und befindet sich in der Beta.
Daher wurde diese Arbeit mit PostgreSQL Version 18.4 durchgeführt~\cite{PostgreSQL}.

Die Architektur von PostgreSQL sieht folgendermaßen aus, es gibt einen Serverprozess der die Datenbank verwaltet,
Verbindungen akzeptiert und die Datenbankaktionen eines Clients ausführt.
Dieses Programm heißt postgres.
Der Client kann dabei vielseitig sein, sei es ein textbasierter Client, eine grafische App oder ein Webserver.
Alle bauen eine TCP/IP-Verbindung zu dem Server auf, welcher mehrere Verbindungen gleichzeitig eingehen kann~\cite{PostgreSQLArchitecture}.

Für PostgreSQL Datenbanken sind Transaktionen das fundamentale Konzept.
Sie bündeln mehrere Schritte in eine einzige Operation.
Die Transaktionen basieren auf dem ACID Prinzip\footnote{Atomicity, eine Transaktion wird entweder ganz oder gar nicht ausgeführt.
	Consistency, nach jeder Transaktion bleiben alle Daten in einem gültigen und korrekten Zustand.
	Isolation, gleichzeitig ablaufende Transaktionen beeinflussen sich nicht gegenseitig.
	Durability, bestätigte Änderungen bleiben dauerhaft gespeichert~\cite{Haerder1983}.}, sodass eine Transaktion immer
gültig oder ungültig ist.
Eine Transaktion beginnt immer mit einem \texttt{BEGIN;} und endet mit einem \texttt{COMMIT;} oder \texttt{ROLLBACK;},
um eine Transaktion zu beenden~\cite{PostgreSQLTransactions}.

Um eine erfolgreiche Recovery bei Abstürzen zu gewährleisten, arbeitet PostgreSQL mit write ahead log (WAL)~\cite{PostgreSQLWAL}.
WAL basiert auf \texttt{UNDO}, Informationen müssen vor der Propagation\footnote{Propagation ist die Operation,
	die geschriebene Seiten zu einem Teil der materialized Database macht. Die materialized Database hingegen ist der
	persistente Datenbankzustand, welcher nach einem Absturz vorhanden ist~\cite{Haerder1983}.}
der entsprechenden Änderung in das Log geschrieben werden, und \texttt{REDO}, Informationen muss in das temp, archive
Log geschrieben werden bevor die Bestätigung des EOT\footnote{End of Transaction} kommt~\cite{Haerder1983}.
